\begin{table}[htbp]
\centering
\caption{Stage~II preference criteria. All criteria are evaluated for the five Stage~I candidates; the four data-derived criteria are computed from open datasets, while the three expert-judged criteria are elicited in the survey.}
\label{tab:s9}
\resizebox{\textwidth}{!}{%
\begin{tabular}{llllp{3.6cm}p{4.2cm}}
\toprule
Code & Criterion & Unit & Direction & Source / mode & Literature support \\
\midrule
K1 & Grid-connection proximity (distance to $\ge$154\,kV substation) & km & Cost & OpenStreetMap / data & Messaoudi et~al. (2019); Ali et~al. (2022); Yaman (2024) \\
K2 & Water stress of the supplying basin & score 0--5 & Cost & World Resources Institute (WRI) Aqueduct 4.0 (Kuzma et~al., 2023) / data & Ishmam et~al. (2024); Tonelli et~al. (2023); Gezen U{\c{c}}ar (2026) \\
K3 & Ecological sensitivity (distance to protected areas beyond the 1\,km buffer) & km & Benefit & OpenStreetMap protected-area snapshot / data & Ali et~al. (2022); Pinto et~al. (2025); Yaman (2024) \\
K4 & Seismic hazard (peak ground acceleration (PGA); 475-year return period) & g & Cost & Global Earthquake Model (GEM) Global Seismic Hazard Model (GSHM) v2023.1 / data & Emre et~al. (2018); Yaman (2024); Ali et~al. (2026); Zhao et~al. (2024) \\
K5 & Local social acceptance & 1--9 & Benefit & Expert survey & Emodi et~al. (2021); Zhao et~al. (2024); Zun and McLellan (2026) \\
K6 & Permitting, land-acquisition and regulatory readiness & 1--9 & Benefit & Expert survey & Pinto et~al. (2025); Zun and McLellan (2026) \\
K7 & Industrial-logistics ecosystem and skilled workforce & 1--9 & Benefit & Expert survey & Ishmam et~al. (2024); Zun and McLellan (2026) \\
\bottomrule
\end{tabular}%
}
\end{table}

\subsection*{Sources for this table}
\begin{enumerate}\setlength{\itemsep}{0pt}
\item Ali, F., Bennui, A., Chowdhury, S., Techato, K., 2022. Suitable site selection for solar-based green hydrogen in southern Thailand using GIS-MCDM approach. Sustainability 14~(11), 6597. \url{https://doi.org/10.3390/su14116597}
\item Ali, T., Aytekin, A., Demidionov, M., 2026. Optimal siting of wind powered green hydrogen production in Bangladesh using interval-valued T-spherical fuzzy multi-criteria decision-making method. International Journal of Hydrogen Energy 217, 153931. \url{https://doi.org/10.1016/j.ijhydene.2026.153931}
\item Emodi, N.V., Lovell, H., Levitt, C., Franklin, E., 2021. A systematic literature review of societal acceptance and stakeholders' perception of hydrogen technologies. International Journal of Hydrogen Energy 46~(60), 30669--30697. \url{https://doi.org/10.1016/j.ijhydene.2021.06.212}
\item Emre, {\"O}., Duman, T.Y., {\"O}zalp, S., {\c{S}}aro{\u{g}}lu, F., Olgun, {\c{S}}., Elmac{\i}, H., et~al., 2018. Active fault database of Turkey. Bulletin of Earthquake Engineering 16~(8), 3229--3275. \url{https://doi.org/10.1007/s10518-016-0041-2}
\item Gezen U{\c{c}}ar, M., 2026. Spatial optimization of green hydrogen production hubs and capacity allocation: A hybrid MCDA--MILP framework for T{\"u}rkiye. International Journal of Hydrogen Energy 269, 157032. \url{https://doi.org/10.1016/j.ijhydene.2026.157032}
\item Ishmam, S., Heinrichs, H., Winkler, C., Bayat, B., Lahnaoui, A., Agbo, S., et~al., 2024. Mapping local green hydrogen cost-potentials by a multidisciplinary approach. International Journal of Hydrogen Energy 87, 1155--1170. \url{https://doi.org/10.1016/j.ijhydene.2024.08.501}
\item Kuzma, S., Bierkens, M.F.P., Lakshman, S., Luo, T., Saccoccia, L., Sutanudjaja, E.H., et~al., 2023. Aqueduct 4.0: Updated decision-relevant global water risk indicators. Technical Note. World Resources Institute, Washington, DC. \url{https://doi.org/10.46830/writn.23.00061}
\item Messaoudi, D., Settou, N., Negrou, B., Settou, B., 2019. GIS based multi-criteria decision making for solar hydrogen production sites selection in Algeria. International Journal of Hydrogen Energy 44~(60), 31808--31831. \url{https://doi.org/10.1016/j.ijhydene.2019.10.099}
\item Pinto, M.C., Gaeta, M., Arco, E., Boccardo, P., Corgnati, S.P., 2025. Mapping the suitability of North Africa for green hydrogen production: an application of a multi-criteria spatial decision support system combining GIS and AHP for Tunisia. Energy, Sustainability and Society 15~(1), 20. \url{https://doi.org/10.1186/s13705-025-00512-4}
\item Tonelli, D., Rosa, L., Gabrielli, P., Caldeira, K., Parente, A., Contino, F., 2023. Global land and water limits to electrolytic hydrogen production using wind and solar resources. Nature Communications 14~(1), 5532. \url{https://doi.org/10.1038/s41467-023-41107-x}
\item Yaman, A., 2024. A GIS-based multi-criteria decision-making approach (GIS-MCDM) for determination of the most appropriate site selection of onshore wind farm in Adana, Turkey. Clean Technologies and Environmental Policy 26~(12), 4231--4254. \url{https://doi.org/10.1007/s10098-024-02866-3}
\item Zhao, H., Ge, Y., Wang, W., 2024. A study on offshore wind farm site selection based on CRITIC and CPT-TOPSIS: a case study of China. Kybernetes 53~(3), 1117--1147. \url{https://doi.org/10.1108/K-09-2022-1267}
\item Zun, M.T., McLellan, B.C., 2026. Where in the world should we produce green hydrogen? An objective first-pass site selection. Hydrogen 7~(1), 11. \url{https://doi.org/10.3390/hydrogen7010011}
\end{enumerate}
